\subsection{Constraints on a symmetric three-plane}
\label{sec:gct-three-plane}

\paragraph{Maximality of the star plane.}
In fact the first-row-and-column level space is maximal by inclusion in
the full symmetric level fiber for every \(n\ge3\). Put \(m=n-1\), let
\(e\) be the all-ones vector of length \(m\), and write a proposed
extension direction as \(B=\bigl(\begin{smallmatrix}d&b^T\\ b&C\end{smallmatrix}\bigr)\).
On adding the star direction with arm \(u\in e^\perp\), the pencil matrix
has top-left entry \(1+td\), arm \(w=e+tb+u\), and lower block
\(M=J_m+tC\). Its permanent is
\begin{equation*}
  (1+td)\operatorname{per}(M)+w^T D(t)w,
  \qquad D_{ij}(t)=\operatorname{per}\bigl(M_{\widehat i,\widehat j}\bigr).
\end{equation*}
The quadratic and linear coefficients in \(u\) require the bilinear
form of \(D(t)\) to vanish on \(e^\perp\times e^\perp\) and between
\(e^\perp\) and \(e+tb\). Here \(D(t)\) is symmetric. Since
\(e^T(e+tb)\ne0\) except possibly at one \(t\), this forces
\(D(t)e^\perp=0\). The identity holds for all \(t\) by polynomiality,
so \(D(t)\) has constant entries for every \(t\).

Let \(r_i\) be the row sums of \(C\) and \(T=\sum_i r_i\).
Differentiation gives
\(D'_{ij}(0)=(m-2)!(T-r_i-r_j+C_{ij})\).
Its row sums are proportional to \(T-r_i\), so constancy forces equal
\(r_i\) and then \(C=cJ_m\). If \(c\ne0\), setting \(t=-1/c\) zeros
\(M\) and its permanental cofactors, making the permanent zero. Thus
\(C=0\). The quadratic coefficient in \(t\) now forces \(e^Tb=0\), and
the linear one forces \(d=0\). The proposed direction belongs to the
original star plane.

\paragraph{Matrix rank.}
Every candidate three-plane at \(J_4\) must contain a direction of
matrix rank at least three. First, a symmetric rank-one direction is
\(aa^T\) for some nonzero \(a\in\mathbb C^4\). The coefficient of \(t^k\)
in \(\operatorname{per}(J_4+taa^T)\) is
\(k!(4-k)!e_k(a)^2\), where \(e_k\) is the elementary symmetric
polynomial. Selecting the \(k\) rows and \(k\) columns contributes
the two factors of \(e_k\). If all four positive-degree coefficients
vanish, the entries of \(a\) are roots of \(z^4\), forcing \(a=0\).

Suppose all matrices in a level three-plane have rank at most two.
By congruence bring one rank-two matrix to \(\operatorname{diag}(I_2,0)\)
and write every other matrix in blocks \(\bigl(\begin{smallmatrix}A&C\\C^T&D\end{smallmatrix}\bigr)\).
The Schur complement of \(I_2+tA\) must vanish in the combined pencil.
Expanding \(tD-t^2C^T(I_2+tA)^{-1}C\) gives
\(D=0\), \(C^TC=0\), and \(C^TAC=0\). If every \(C\) is zero, the
three-plane is all symmetric \(2\times2\) matrices in the top block
and contains a rank-one matrix. Otherwise the linear space of \(C\)'s
has image in one common isotropic line \(\mathbb Cp\) in \(\mathbb C^2\).
Indeed, \(C^TC=0\) places each image in one of two isotropic lines,
and a linear space contained in the union of two subspaces lies in one.
The equality \(C^TAC=0\) forces \(p^TAp=0\) whenever \(C\ne0\), hence
for every matrix by linearity. These \(A\)'s have the form
\(px^T+xp^T\), since this two-dimensional image lies in the kernel of
the nonzero linear functional \(A\mapsto p^TAp\). Undoing the
congruence, every direction is therefore \(av^T+va^T\) for a common
vector \(a\ne0\) and \(v\) in a three-dimensional space \(U_0\). The
map from \(v\) is injective by considering a coordinate where \(a_i\ne0\).

Write \(K=J_4-I_4\) and \(q(v)=v^TKv\). The linear coefficient forces
\((\sum a_i)(\sum v_i)=0\), and the Hessian forces
\begin{equation*}
  q(a)q(v)+(a^TKv)^2=0\qquad(v\in U_0).
\end{equation*}
If \(\sum a_i\ne0\), then \(U_0=\mathbf1^\perp\), where \(q\) has
rank three. The equation forces \(q(a)=0\) by comparison with a form
of rank at most one. It then forces \(a\) to be proportional to
\(\mathbf1\), contradicting \(q(a)=0\). If \(\sum a_i=0\) and
\(q(a)\ne0\), the displayed form on \(\mathbb C^4\) is nondegenerate.
In fact a kernel vector satisfies \(q(a)v+a(a^TKv)=0\), which first
gives \(a^TKv=0\), then \(v=0\). A three-dimensional totally isotropic
space is impossible in this four-dimensional form. If \(q(a)=0\),
then \(U_0=\{v:a^TKv=0\}\) contains \(a\), producing the prohibited
rank-one matrix \(2aa^T\).

\paragraph{A finite symmetry obstruction.}
Let the Klein four group act on the four indices by its regular
permutations, simultaneously on rows and columns. A stable level space
decomposes by the four character projectors. For a nontrivial character
\(\chi\), the symmetric matrix space of that character has basis
\(\mathbf1\odot\chi,\ \psi\odot\phi\), where \(x\odot y=xy^T+yx^T\)
and \(\psi,\phi\) are the other nontrivial characters. For coefficients
\(u,v\), the quadratic and quartic permanent coefficients are
\(32(v^2-3u^2)\) and \(64(u^2+v^2)^2\). They force \(u=v=0\).
The trivial-character matrices have first row \((d,a,b,c)\) with the
others given by translations. The linear coefficient forces
\(d=-a-b-c\), and the cubic coefficient is \(-8(a+b)(a+c)(b+c)\).
Relabeling the nonzero group elements allows \(b=-a\). The quadratic
coefficient is then \(8(a^2+c^2)\) and the quartic reduces to
\(-8a^4\) modulo \(a^2+c^2\), forcing the zero direction. Thus the
only stable level space for this simultaneous Klein four action is
zero. The exact coefficient checks in \path{gct_permanent_fiber.py}
cover all four characters.

\paragraph{Coefficient charts.}
Paper~I of the Jacobian collision project isolates the residual factor
in the generic collision algebra of a cubic Keller map~\cite{jc3collision}.
Here the corresponding reduction is to impose the quadratic coefficient
identically, then examine the cubic and quartic coefficients on that locus.
Write \(q_j(B)=[t^j]\operatorname{per}(J_4+tB)\). In the orthonormal
Hadamard basis consisting of half the all-ones vector and half the
three nontrivial Klein characters, write the tangent matrix as
\(\left(\begin{smallmatrix}0&b^T\\b&C\end{smallmatrix}\right)\).
For \(\{i,j,k\}=\{1,2,3\}\), use the coordinate vectors \(d_i,c_i,b_i\)
for \(C_{ii}, C_{jk}=C_{kj}\), and the \(i\)-th arm entry. The Hessian
formula gives \(q_2=\sum_i(C_{ii}^2+2C_{jk}^2-6b_i^2)\), with each
opposite entry \(C_{jk}\) counted once. For its polar form
\(\beta(u,v)=(q_2(u+v)-q_2(u)-q_2(v))/2\), a Witt basis over
\(\mathbb Q(\mathrm i)\) is
\begin{align*}
 e_i &= 2d_i+c_i+b_i,& f_i &= (2d_i-c_i-b_i)/8,& w_i &=3c_i+b_i
    && (1\le i\le 3),\\
 e_4 &= w_1+\mathrm i w_2,& f_4 &= (w_1-\mathrm i w_2)/24,& h &= w_3.
\end{align*}
The four pairs have \(\beta(e_i,f_i)=1\) and zero self-pairings.
Distinct pairs are orthogonal, and \(h\) is orthogonal to them with
\(q_2(h)=12\).

Every totally isotropic three-plane has a nonzero projection minor
onto three coordinate lines chosen from distinct pairs. To see this,
scale the two lines of each pair by opposite powers of a parameter
\(\lambda\), with weights \(1,2,3,4\) on the \(e_i\), and fix \(h\).
As \(\lambda\to0\), the lowest-weight terms in the plane's Pl\"ucker
coordinates give a limit point of the closed isotropic Grassmannian.
The limit is fixed under this scaling, hence is a coordinate plane
since all nine weights are distinct. Isotropy excludes \(h\) and excludes
selecting both lines of any pair. A nonzero minor of the limit was
already nonzero on the original plane. This gives 32 projection charts.

On each chart, split the coordinates into the selected lines, their
dual lines, and the remaining three-dimensional orthogonal complement.
The normalized frame is \((I_3,A,M)^T\) in these blocks. With \(Q\) the
Gram matrix on the complement, isotropy is equivalent to
\(A+A^T+M^TQM=0\). Thus \(A=S-M^TQM/2\), where \(S\) is skew-symmetric,
parametrizes the chart by 12 variables. For the frame \(B_1,B_2,B_3\),
the residual equations are the ten coefficients of
\(q_3(\sum_i t_i B_i)\) and the fifteen of \(q_4(\sum_i t_i B_i)\).

For the existence question it suffices to use the four pivot triples
\((e_2,e_3,e_4)\), \((e_1,e_2,e_4)\), \((e_1,e_2,e_3)\), and
\((f_1,e_2,e_3)\). Klein translations change the signs of \(c_i,b_i\)
for an even set of \(i\le3\), interchanging the corresponding \(e_i,f_i\)
up to scale. They preserve or swap the fourth pair up to scale.
Complex conjugation independently swaps the fourth pair. If a pair
among the first three is omitted, these operations allow all orientations
on the selected pairs. The index permutation interchanging the first two
characters also identifies omissions one and two. If the fourth pair
is omitted, even swaps leave two orientation parity classes.
All these operations preserve the permanent coefficient locus.
A common zero on any representative chart constructs a three-plane.
Unit ideals on all four would certify emptiness over \(\mathbb C\).
The executable checks the basis and chart symmetries and generates
the residual systems.

\paragraph{Excluding planes with zero row sums.}
If every direction has zero row sums, the arm \(b\) is zero. Let \(H\)
denote the three nontrivial columns of the normalized Hadamard matrix,
so \(B=HCH^T\). On this slice, \(q_2(B)=\operatorname{tr}(C^2)\) and
\(q_3(B)=\operatorname{per}(C)\). Indeed, \(q_3\) is one sixth the sum of
products of three entries in distinct rows and columns. The indicator
of distinct rows is one minus the sum of the three pairwise equality
indicators plus twice the all-equal indicator, and likewise for columns.
A term with a singleton row or column index sums to zero. This leaves
\(q_3(B)=\tfrac23\sum_{ij}B_{ij}^3\). The sum of a product of three
columns of \(H\) over the four indices is \(1/2\) for three distinct
characters and zero otherwise. Hence \(\sum_{ij}B_{ij}^3=\tfrac32
\operatorname{per}(C)\).

Write \(x_i=C_{ii}\) and \(y_i=C_{jk}\) for opposite \(j,k\). Fix
\(\kappa^2=-1/2\). The same projection-minor argument gives eight charts
for a totally isotropic three-plane in \(\sum_i(x_i^2+2y_i^2)=0\).
For signs \(\epsilon_i=\pm1\), set \(y_i=\kappa\epsilon_i z_i\).
The pivot projection is onto \((x+z)/2\), so the chart is
\(x=(I_3+S)t\), \(z=(I_3-S)t\) for a skew-symmetric \(S\) with three
entries above the diagonal. Then
\begin{equation*}
 2\operatorname{per}(C)=2x_1x_2x_3-\sum_i x_i z_i^2
             -2\kappa(\epsilon_1\epsilon_2\epsilon_3)z_1z_2z_3.
\end{equation*}
It suffices to check sign product \(+1\) by conjugation. Take the ten
coefficients in \(t\) of the displayed cubic, and multiply each by
the monomials in the three skew entries that keep the product at
degree at most four. In the basis of all 35 monomials of degree at most
four, the resulting \(46\times35\) coefficient matrix has rank 35
under reduction to \(\mathbb F_{11}\) with \(\kappa\mapsto4\). The
executable checks this rank. The reduction respects \(2\kappa^2=-1\),
so the nonzero minor also proves rank 35 over \(\mathbb Q(\kappa)\).
In particular the row span contains the constant one, making the cubic
coefficient ideal the unit ideal. Thus any level three-plane must
contain a direction with nonzero row sums. The unrestricted four chart
systems still require a common zero or an emptiness certificate.

\paragraph{A three-cycle and its eigenspaces.}
A simultaneous index permutation fixing index zero permutes the three
Klein characters and thus permutes each of the three coordinate
triples \(b,c,d\). A three-cycle has one fixed sector and two nontrivial
eigenspaces. In the fixed sector the three triples are each constant,
with respective entries \(b,c,d\). The coefficient ideal
\((q_2,q_3,q_4)\) is the unit ideal on the charts \(b=1\) and
\(b=0,c=1\), as checked over \(\mathbb Q\), \(\mathbb F_5\), and
\(\mathbb F_7\). The remaining projective direction with \(b=c=0\)
has \(q_2=3d^2\). Thus no nonzero fixed direction is level.

In a nontrivial eigenspace the three triples are scalar multiples of
\((1,\omega,\omega^2)\), or of its conjugate, with
\(\omega^2+\omega+1=0\). Invariance forces \(q_2=q_4=0\) there.
The cubic coefficient with \(u=b+c\) is
\begin{equation*}
 F(b,c,d)=d^3+3u^2d+8c^3-12cu^2+6u^3.
\end{equation*}
A singular point of this cubic would satisfy
\(d^2+u^2=0\), \(2c^2-u^2=0\) and \(u(d-4c+3u)=0\).
If \(u=0\) these force all coordinates to vanish. If \(u\ne0\)
normalize \(u=1\). Then \(d=4c-3\) and \(c^2=1/2\), so
\(d^2=-1\) forces \(c=3/4\), which contradicts \(c^2=1/2\).
This also holds in characteristics 5 and 7. A smooth projective
plane cubic contains no line, since a linear component would meet
the residual conic at a singular point. Hence a level subspace in
either nontrivial eigenspace has dimension at most one.
A level subspace invariant under a three-cycle therefore has
dimension at most two.

The simultaneous \(S_3\) fixing index zero has only one invariant
level two-plane over \(\mathbb C\). Indeed the tangent representation
is three copies of the three-point permutation representation. A
stable level two-plane has no fixed vectors and is a standard copy.
Consequently its triples are \(br,cr,dr\) for a common vector \(r\)
with \(r_1+r_2+r_3=0\) and scalar profile \((b,c,d)\). Let
\(N=r_1^2+r_1r_2+r_2^2\). The exact coefficient identities are
\begin{equation*}
 q_2=2(d^2+2c^2-6b^2)N,\qquad
 q_3=r_1r_2r_3 F,\qquad
 4q_4=((b-c)^2-d^2)^2 N^2.
\end{equation*}
The last identity gives \(d=\pm(b-c)\). The first then gives
\((c+b)(3c-5b)=0\), with \(b\ne0\) for a nonzero profile.
Normalizing \(b=1\), the cubic vanishes only at \((c,d)=(-1,2)\).
This gives the star plane at index zero.

For any level three-plane, its stabilizer in the simultaneous
\(S_4\) contains no element of order three. A subgroup of order
eight would contain the regular Klein four, which is also excluded.
Thus the stabilizer has order one, two, or four. Every candidate
has an orbit of 24, 12, or six distinct relabelings. In particular,
a finite level three-plane locus has a number of geometric points
divisible by six.
